\documentclass[a4paper,12pt]{article}
\usepackage{graphicx}
\usepackage{caption}
\usepackage{subcaption}
\usepackage{float}
\usepackage[margin=1in]{geometry}
\usepackage{lmodern}
\usepackage{amsmath}
\usepackage{xcolor}
\usepackage{listings}
\usepackage{pgffor}

\lstset{
  breaklines=true,
  postbreak=\mbox{\textcolor{red}{$\hookrightarrow$}\space},
}

\graphicspath{{../../python/kNN/output/}}

\begin{document}

\title{Relatório kNN}
\author{Eduardo}
\date{\today}
\maketitle

\section*{Exercício 1}

As figuras a seguir relacionam as superfícies de decisão e os espaços característicos aos parametros avaliados.
  
\foreach \k in {1, 25, 50} {
  \foreach \r in {1.00, 2.00, 3.00} {
    \begin{figure}[H]
      \centering
      \includegraphics[width=0.5\textwidth]{kNN_\k_noise\r.png}
      \caption{Superfície de decisão e espaço característico para k=\k\ e r=\r}
    \end{figure}
  }
}

Pelos resultados, fica claro que, quanto maior a quantidade de vizinhos avaliada, maior a robustês do classificador em relação a \textit{outliers}. No entanto, valores muito altos podem fazer com que este não capture com fidelidade a distribuição espacial dos dados.

\section*{Exercício 2}

As figuras a seguir relacionam as superfícies de decisão e os espaços característicos aos parametros avaliados.

\foreach \k in {1, 25, 50} {
  \foreach \h in {0.00, 0.01, 1.00} {
    \begin{figure}[H]
      \centering
      \includegraphics[width=0.5\textwidth]{kNN_\k_h\h.png}
      \caption{Superfície de decisão e espaço característico para k=\k\ e h=\h}
    \end{figure}
  }
}

Os resultados indicam que a combinação de \texttt{k} e \texttt{h} é diretamente refletida no espaço característico. Nota-se, também, que a distribuição das amostras neste é fortemente relacionada à superfície de decisão no espaço de entrada.

\end{document}
